\begin{titlepage}
	\definecolor{forestgreen}{RGB}{34,92,56}   % 森林绿主题色
	\begin{tikzpicture}[remember picture, overlay]
		% ---- 1. 墨绿渐变铺满整个封面 ----
		\shade[top color=forestgreen, middle color=forestgreen!90, bottom color=forestgreen!60!black]
			(current page.south west) rectangle (current page.north east);

		% ---- 2. 细点阵（模拟网格，微弱纹理）----
		\foreach \x in {1,...,20} {
			\foreach \y in {1,...,29} {
				\fill[white, opacity=0.05] (\x,\y) circle[radius=0.018cm];
			}
		}

		% ---- 3. 主标题后方的柔光晕 ----
		\fill[white, opacity=0.05] ([yshift=-8.2cm]current page.north) circle[radius=7.0cm];
		\fill[white, opacity=0.05] ([yshift=-8.2cm]current page.north) circle[radius=4.6cm];

		% ---- 4. 右上：光滑曲面补片（二维流形）与坐标曲线 ----
		\begin{scope}[shift={([xshift=-6.2cm,yshift=-4.7cm]current page.north east)}]
			% 坐标曲线族（水平向）
			\foreach \k in {0,...,5} {
				\draw[white, opacity=0.55, line width=0.55pt, domain=0:5.0, samples=60, smooth]
					plot (\x, {0.2 + 0.28*\k + 0.18*sin(deg(1.2*\x + 0.5*\k))});
			}
			% 坐标曲线族（垂直向）
			\foreach \k in {0,...,6} {
				\draw[white, opacity=0.55, line width=0.55pt, domain=0:1.6, samples=60, smooth]
					plot ({0.3 + 0.72*\k + 0.20*sin(deg(1.8*\x + 0.4*\k))}, \x);
			}
			% 曲面边界
			\draw[white, line width=1.0pt, opacity=0.85, domain=0:5.0, samples=80, smooth]
				plot (\x, {1.60 + 0.20*sin(deg(1.2*\x))});
			\draw[white, line width=1.0pt, opacity=0.85, domain=0:5.0, samples=80, smooth]
				plot (\x, {0.20 + 0.20*sin(deg(1.2*\x + 0.7))});
			% 局部坐标卡 U（青色高亮）
			\fill[cyan!65!white, opacity=0.18] (2.65,0.90) circle[radius=0.85cm];
			\draw[cyan!65!white, dashed, line width=0.6pt] (2.65,0.90) circle[radius=0.85cm];
			\node[anchor=west, text=cyan!65!white, font=\footnotesize\itshape] at (3.05,1.40) {$U$};
		\end{scope}

		% ---- 5. 左下：球面上的局部坐标卡（球面 + 经纬网格 + 开邻域 U）----
		\begin{scope}[shift={([xshift=0.9cm,yshift=0.9cm]current page.south west)}]
			% 球面轮廓
			\draw[white, line width=1.0pt, opacity=0.85] (1.15,1.25) circle (1.1cm);
			% 经纬网格
			\draw[white, dashed, opacity=0.40, line width=0.5pt] (1.15,1.25) ellipse (1.1cm and 0.40cm);
			\draw[white, dashed, opacity=0.40, line width=0.5pt, rotate around={45:(1.15,1.25)}] (1.15,1.25) ellipse (1.1cm and 0.40cm);
			\draw[white, dashed, opacity=0.40, line width=0.5pt, rotate around={-45:(1.15,1.25)}] (1.15,1.25) ellipse (1.1cm and 0.40cm);
			\draw[white, dashed, opacity=0.40, line width=0.5pt, rotate around={90:(1.15,1.25)}] (1.15,1.25) ellipse (1.1cm and 0.40cm);
			% 局部坐标卡 U（青色高亮开邻域）
			\fill[cyan!65!white, opacity=0.25] (1.68,1.68) circle[radius=0.40cm];
			\draw[cyan!65!white, dashed, line width=0.6pt] (1.68,1.68) circle[radius=0.40cm];
			% U 上的局部坐标轴
			\draw[cyan!65!white, line width=0.9pt, opacity=0.95, -{Stealth}] (1.68,1.68) -- (2.00,1.50);
			\draw[cyan!65!white, line width=0.9pt, opacity=0.95, -{Stealth}] (1.68,1.68) -- (1.84,1.98);
			\node[anchor=west, text=cyan!65!white, font=\footnotesize\itshape] at (2.00,1.88) {$U$};
			\node[anchor=north, text=white, opacity=0.75, font=\footnotesize\itshape] at (1.15,2.42) {$M$};
		\end{scope}

		% ---- 6. 金色细边框 ----
		\draw[gold, line width=0.7pt, opacity=0.55]
			([shift={(0.85,0.85)}]current page.south west) rectangle
			([shift={(-0.85,-0.85)}]current page.north east);

		% ---- 7. 顶部眉题（英文小字 + 金色短线）----
		\coordinate (T) at ([yshift=-3.5cm]current page.north);
		\draw[gold, line width=1.0pt] (T) ++(-3.4cm,0) -- ++(1.0cm,0);
		\draw[gold, line width=1.0pt] (T) ++(2.4cm,0) -- ++(1.0cm,0);
		\node[anchor=north, text=white, opacity=0.85, align=center] at (T) {
			{\sffamily\fontsize{13}{16}\selectfont
				\uppercase{Differential \ Manifolds}}
		};

		% ---- 8. 主标题（居中，使用 \notename）----
		\node[anchor=center, align=center, text=white] at ([yshift=-8.0cm]current page.north) {%
			{\fontsize{40}{50}\sffamily\bfseries \notename}
		};

		% ---- 9. 金色双细线 + 小球面纹饰（紧致流形符号）----
		\coordinate (C) at ([yshift=-11.4cm]current page.north);
		\draw[gold, line width=1.1pt] (C) ++(-2.2cm,0) -- ++(1.55cm,0);
		\draw[gold, line width=1.1pt] (C) ++(0.65cm,0) -- ++(1.55cm,0);
		\draw[gold, line width=0.9pt] (C) circle[radius=0.16cm];
		\draw[gold, line width=0.6pt, opacity=0.75] (C) ellipse (0.16cm and 0.06cm);
		\draw[gold, line width=0.6pt, opacity=0.75, rotate around={60:(C)}] (C) ellipse (0.16cm and 0.06cm);
		\draw[gold, line width=0.6pt, opacity=0.75, rotate around={-60:(C)}] (C) ellipse (0.16cm and 0.06cm);

		% ---- 10. 副标题 ----
		\node[anchor=north, text=white, opacity=0.9, align=center]
			at ([yshift=-12.2cm]current page.north) {
			{\fontsize{16}{20}\sffamily 学\ \ 习\ \ 笔\ \ 记}
		};

		% ---- 11. 底部作者、日期信息 ----
		\coordinate (B) at ([yshift=3.4cm]current page.south);
		\draw[gold, line width=0.8pt, opacity=0.7] (B) ++(-1.1cm,0.85cm) -- ++(2.2cm,0);
		\node[anchor=south, text=white, align=center] at (B) {
			\begin{tabular}{c}
				{\fontsize{20}{24}\sffamily\bfseries \noteauthor} \\[0.4cm]
				{\large \sffamily \today}
			\end{tabular}
		};
	\end{tikzpicture}
\end{titlepage}
